% ======================================================================
% 3.2 Высшие производные и пределы прогноза. Ответственный: B — Попадюк Евгений.
%
% Что сдать (main.tex: теория 1077–1120, порядок 1121–1170,
% разобранный пример 1172–1276; требования к отчёту 1871–1879):
%   1. Вторые частные производные T(lambda,c); равенство смешанных;
%      определитель матрицы Гессе и объяснение, почему он равен нулю.
%   2. Доказать по индукции формулу n-й производной по lambda;
%      записать ряд Тейлора T(lambda+h), найти его сумму и радиус сходимости.
%   3. При росте трафика h=s/2 вычислить прогнозы Тейлора 1, 2, 3 порядков
%      и их ошибки; сделать вывод.
%   4. По формуле Тейлора 2-го порядка для двух переменных спрогнозировать
%      задержку при (Delta lambda, Delta c) своего варианта; сравнить с точной.
%   5* (по желанию, на оценку не влияет): функция выгоды G, стационарная точка,
%      критерий Сильвестра; объяснить результат.
%
% Параметры варианта 3 (приложение Б): mu=80, lambda=600, c=8;
% (Delta lambda, Delta c)=(20, 0.1); h=s/2; для п. 5*: a=1, b=1, p=0.03.
% Рабочая точка и s — из п. 3.1 (файл 31-equilibrium.tex).
%
% Программные значения для таблиц — task_3_2 (python experiments.py 3.2);
% графики ошибок прогнозов от h — п. 3.5.2 (твой же файл 352-taylor-errors.tex).
% ======================================================================

\subsection{Высшие производные и пределы прогноза}

\taskhead{Исходные данные}
\todo{B: рабочая точка варианта 3 и запас $s$ (из п. 3.1); приращения $(\Delta\lambda,\Delta c)=(20,\,0.1)$; величина $h=s/2$.}

\taskhead{Матрица Гессе}
\todo{B: найти вторые частные производные $T(\lambda,c)$, проверить равенство смешанных; составить матрицу Гессе, вычислить определитель и объяснить, почему он равен нулю (связь с п. 3.1, направление нулевого изменения).}

\taskhead{Ряд Тейлора по трафику}
\todo{B: доказать по индукции формулу $\partial^nT/\partial\lambda^n$; записать ряд Тейлора $T(\lambda+h)$ по степеням $h$, найти его сумму и радиус сходимости (пример: main.tex, строки 1182–1192).}

\taskhead{Прогнозы при росте трафика на $h=s/2$}
\todo{B: вычислить прогнозы $P_1,P_2,P_3$, точное значение задержки и ошибки прогнозов; сделать вывод (пример: main.tex, строки 1193–1208).}

\taskhead{Прогноз по двум переменным}
\todo{B: по формуле Тейлора 2-го порядка спрогнозировать задержку при $(\Delta\lambda,\Delta c)=(20,\,0.1)$; сравнить с точным значением (пример: main.tex, строки 1209–1240).}

\taskhead{Сопоставление с программным расчётом}
% Программные значения — вызовы своих функций из solution.py (task_3_2).
\todo{B: листинг вызовов функций и запуск python experiments.py 3.2.}
\begin{table}[htbp]
\centering\footnotesize
\setlength{\tabcolsep}{4pt}
\caption{Сравнение аналитических и программных значений (п. 3.2)}
\label{tab:3-2}
\begin{tabular}{|l|r|r|r|r|}\hline
Величина & \shortstack{Аналитическое\\значение} & \shortstack{Программное\\значение} & \shortstack{Абсолютная\\ошибка} & \shortstack{Относительная\\ошибка}\\\hline
$P_1$ при $h=s/2$ & — & — & — & —\\
$P_2$ при $h=s/2$ & — & — & — & —\\
$P_3$ при $h=s/2$ & — & — & — & —\\
$P_2$ по двум переменным & — & — & — & —\\
$T$ точное & — & — & — & —\\
\hline\end{tabular}
\end{table}

\taskhead{Вывод}
\todo{B: как порядок прогноза влияет на ошибку; где прогноз перестаёт работать.}

% Пункт 3.2.5* (функция выгоды G, стационарная точка, критерий Сильвестра) —
% по желанию, на оценку не влияет. Если выполняется: производные G, стационарная
% точка, матрица Гессе в ней, обоснованный вывод о характере точки.
